\textbf{Phần 1. Trắc nghiệm lựa chọn}
\setcounterpageref{ex}{0}

\begin{ex} %[5H5N28]
Bộ đồ chơi xếp hình Tangram truyền thống gồm có bao nhiêu miếng ghép?
\choice
{$5$ miếng ghép}
{$6$ miếng ghép}
{$8$ miếng ghép}
{\True $7$ miếng ghép}
\loigiai{
Bộ đồ chơi xếp hình Tangram gồm có $7$ miếng ghép (gồm $5$ hình tam giác, $1$ hình vuông và $1$ hình bình hành).\\
\textbf{Đáp án của câu hỏi là D}
}
\end{ex}

\begin{ex} %[5H5N28]
Trong bộ xếp hình Tangram $7$ miếng ghép, có bao nhiêu miếng ghép là hình tam giác?
\choice
{$3$ miếng}
{$4$ miếng}
{\True $5$ miếng}
{$6$ miếng}
\loigiai{
Trong bộ $7$ miếng ghép gồm có $5$ hình tam giác vuông cân, $1$ hình vuông và $1$ hình bình hành.\\
\textbf{Đáp án của câu hỏi là C}
}
\end{ex}

\begin{ex} %[5H5H28]
Để tạo hình mái nhà trong bức tranh ngôi nhà, người ta thường dùng hình phẳng nào dưới đây?
\choice
{Hình tròn}
{\True Hình thang cân hoặc hình tam giác}
{Hình chữ nhật}
{Hình thoi}
\loigiai{
Mái nhà thường có dạng hình thang cân hoặc hình tam giác.\\
\textbf{Đáp án của câu hỏi là B}
}
\end{ex}

\begin{ex} %[5H5H28]
Một bông hoa được tạo thành từ $1$ hình tròn màu vàng làm nhụy và $5$ hình tròn nhỏ màu đỏ có cùng bán kính $1\text{ cm}$ làm cánh hoa. Bán kính hình tròn nhụy hoa là $2\text{ cm}$. Diện tích của nhụy hoa đó là
\choice
{$6{,}28\text{ cm}^2$}
{$3{,}14\text{ cm}^2$}
{$15{,}7\text{ cm}^2$}
{\True $12{,}56\text{ cm}^2$}
\loigiai{
Diện tích nhụy hoa là: $3{,}14 \times 2 \times 2 = 12{,}56\text{ (cm}^2\text{)}$.\\
\textbf{Đáp án của câu hỏi là D}
}
\end{ex}

\begin{ex} %[5H5H28]
Một mái nhà có dạng hình thang với hai đáy là $8\text{ m}$ và $12\text{ m}$, chiều cao là $3\text{ m}$. Diện tích mái nhà đó là
\choice
{$60\text{ m}^2$}
{\True $30\text{ m}^2$}
{$20\text{ m}^2$}
{$40\text{ m}^2$}
\loigiai{
Diện tích mái nhà là: $\dfrac{(8 + 12) \times 3}{2} = 30\text{ (m}^2\text{)}$.\\
\textbf{Đáp án của câu hỏi là B}
}
\end{ex}

\begin{ex} %[5H5H28]
Một bức tường hình chữ nhật có kích thước $6\text{ m} \times 4\text{ m}$. Trên tường có $1$ cửa ra vào hình chữ nhật $1\text{ m} \times 2\text{ m}$ và $1$ cửa sổ hình vuông cạnh $1\text{ m}$. Diện tích phần tường cần sơn là
\choice
{\True $21\text{ m}^2$}
{$24\text{ m}^2$}
{$22\text{ m}^2$}
{$20\text{ m}^2$}
\loigiai{
Diện tích bức tường là: $6 \times 4 = 24\text{ (m}^2\text{)}$.\\
Diện tích các cửa là: $1 \times 2 + 1 \times 1 = 3\text{ (m}^2\text{)}$.\\
Diện tích phần tường cần sơn là: $24 - 3 = 21\text{ (m}^2\text{)}$.\\
\textbf{Đáp án của câu hỏi là A}
}
\end{ex}

\begin{ex} %[5H5H28]
Ghép hai hình tam giác vuông bằng nhau có các cạnh góc vuông là $3\text{ cm}$ và $4\text{ cm}$ thành một hình chữ nhật. Chu vi hình chữ nhật đó là
\choice
{$7\text{ cm}$}
{$12\text{ cm}$}
{$24\text{ cm}$}
{\True $14\text{ cm}$}
\loigiai{
Hình chữ nhật tạo thành có chiều dài $4\text{ cm}$ và chiều rộng $3\text{ cm}$.\\
Chu vi là: $(4 + 3) \times 2 = 14\text{ (cm)}$.\\
\textbf{Đáp án của câu hỏi là D}
}
\end{ex}

\begin{ex} %[5H5V28]
Từ một tờ giấy hình vuông có diện tích $64\text{ cm}^2$, bạn Mai cắt ra thành $7$ miếng ghép Tangram rồi ghép lại thành một hình con thỏ. Diện tích hình con thỏ đó bằng
\choice
{\True $64\text{ cm}^2$}
{$32\text{ cm}^2$}
{$128\text{ cm}^2$}
{$16\text{ cm}^2$}
\loigiai{
Vì hình con thỏ được ghép lại từ toàn bộ $7$ mảnh của hình vuông ban đầu nên diện tích không đổi và bằng $64\text{ cm}^2$.\\
\textbf{Đáp án của câu hỏi là A}
}
\end{ex}

\begin{ex} %[5H5V28]
Một cánh diều có dạng hình thoi gồm hai đường chéo dài $60\text{ cm}$ và $80\text{ cm}$. Diện tích của cánh diều đó là
\choice
{$4\,800\text{ cm}^2$}
{$1\,200\text{ cm}^2$}
{\True $2\,400\text{ cm}^2$}
{$240\text{ cm}^2$}
\loigiai{
Diện tích hình thoi bằng nửa tích hai đường chéo: $\dfrac{60 \times 80}{2} = 2\,400\text{ (cm}^2\text{)}$.\\
\textbf{Đáp án của câu hỏi là C}
}
\end{ex}

\begin{ex} %[5H5V28]
Để làm một chiếc thuyền buồm origami, bạn Việt dùng một tờ giấy hình chữ nhật có kích thước $15\text{ cm} \times 20\text{ cm}$. Diện tích của tờ giấy đó là
\choice
{\True $300\text{ cm}^2$}
{$35\text{ cm}^2$}
{$70\text{ cm}^2$}
{$150\text{ cm}^2$}
\loigiai{
Diện tích tờ giấy là: $15 \times 20 = 300\text{ (cm}^2\text{)}$.\\
\textbf{Đáp án của câu hỏi là A}
}
\end{ex}

\begin{ex} %[5H5V28]
Bạn Nam cắt một tấm bìa hình vuông cạnh $20\text{ cm}$ thành $4$ hình tam giác bằng nhau bằng hai đường chéo. Diện tích của mỗi hình tam giác đó là
\choice
{$50\text{ cm}^2$}
{\True $100\text{ cm}^2$}
{$200\text{ cm}^2$}
{$400\text{ cm}^2$}
\loigiai{
Diện tích hình vuông ban đầu là: $20 \times 20 = 400\text{ (cm}^2\text{)}$.\\
Diện tích mỗi hình tam giác là: $400 : 4 = 100\text{ (cm}^2\text{)}$.\\
\textbf{Đáp án của câu hỏi là B}
}
\end{ex}

\begin{ex} %[5H5V28]
Người ta cắt $4$ góc của một tấm bìa hình vuông cạnh $30\text{ cm}$, mỗi góc bị cắt là một hình vuông nhỏ cạnh $5\text{ cm}$. Diện tích phần bìa còn lại là
\choice
{$900\text{ cm}^2$}
{$850\text{ cm}^2$}
{\True $800\text{ cm}^2$}
{$700\text{ cm}^2$}
\loigiai{
Diện tích tấm bìa ban đầu: $30 \times 30 = 900\text{ (cm}^2\text{)}$.\\
Diện tích $4$ hình vuông nhỏ bị cắt: $4 \times (5 \times 5) = 100\text{ (cm}^2\text{)}$.\\
Diện tích phần bìa còn lại: $900 - 100 = 800\text{ (cm}^2\text{)}$.\\
\textbf{Đáp án của câu hỏi là C}
}
\end{ex}

\textbf{Phần 2. Trắc nghiệm đúng sai}
\setcounterpageref{ex}{0}

\begin{ex} %[5H5H28]
Xét tính đúng sai của các khẳng định sau về các phép đo, vẽ, ghép hình phẳng:
\choiceTF
{\True Khi ghép hai tam giác vuông cân bằng nhau ta có thể tạo thành một hình vuông.}
{Khi cắt một hình phẳng thành các mảnh nhỏ rồi ghép lại thì tổng diện tích bị giảm đi.}
{\True Bộ ghép hình Tangram gồm đúng $5$ hình tam giác, $1$ hình vuông và $1$ hình bình hành.}
{Gấp đôi một tờ giấy hình vuông theo đường chéo ta nhận được một hình thang vuông.}
\loigiai{
\begin{itemize}
    \item Hai tam giác vuông cân bằng nhau ghép theo cạnh huyền tạo thành một hình vuông.
    \item Diện tích hình được bảo toàn khi cắt ghép, không bị giảm đi.
    \item Tangram gồm 7 miếng: 5 tam giác, 1 hình vuông, 1 hình bình hành.
    \item Gấp đôi hình vuông theo đường chéo ta được tam giác vuông cân, không phải hình thang.
\end{itemize}
\begin{itemize}
    \item a) \textbf{Đúng}.
    \item b) \textbf{Sai}.
    \item c) \textbf{Đúng}.
    \item d) \textbf{Sai}.
\end{itemize}
}
\end{ex}

\begin{ex} %[5H5V28]
Một bức tranh trang trí gồm một hình vuông lớn cạnh $10\text{ cm}$ và một hình tròn nội tiếp bên trong tiếp xúc với bốn cạnh của hình vuông (đường kính $10\text{ cm}$). Xét tính đúng sai của các khẳng định sau:
\choiceTF
{\True Diện tích của hình vuông lớn là $100\text{ cm}^2$.}
{\True Bán kính của hình tròn là $5\text{ cm}$.}
{\True Diện tích của hình tròn là $78{,}5\text{ cm}^2$.}
{Diện tích phần trang trí còn lại (bên ngoài hình tròn, bên trong hình vuông) là $22{,}5\text{ cm}^2$.}
\loigiai{
\begin{itemize}
    \item Diện tích hình vuông: $10 \times 10 = 100\text{ cm}^2$.
    \item Bán kính hình tròn: $10 : 2 = 5\text{ cm}$.
    \item Diện tích hình tròn: $3{,}14 \times 5 \times 5 = 78{,}5\text{ cm}^2$.
    \item Diện tích phần còn lại: $100 - 78{,}5 = 21{,}5\text{ cm}^2 \neq 22{,}5\text{ cm}^2$.
\end{itemize}
\begin{itemize}
    \item a) \textbf{Đúng}.
    \item b) \textbf{Đúng}.
    \item c) \textbf{Đúng}.
    \item d) \textbf{Sai}: Kết quả là $21{,}5\text{ cm}^2$.
\end{itemize}
}
\end{ex}

\textbf{Phần 3. Trả lời ngắn}
\setcounterpageref{ex}{0}

\begin{ex} %[5H5H28]
Một mái nhà hình thang có hai đáy là $8\text{ m}$ và $12\text{ m}$, chiều cao $3\text{ m}$. Tính diện tích mái nhà đó theo mét vuông.
\shortans{30}
\loigiai{
\textbf{Đáp số:} $30$\\
$\dfrac{(8 + 12) \times 3}{2} = 30\text{ (m}^2\text{)}$.
}
\end{ex}

\begin{ex} %[5H5H28]
Tờ giấy hình chữ nhật có kích thước $15\text{ cm} \times 20\text{ cm}$. Tính diện tích tờ giấy đó theo xăng-ti-mét vuông.
\shortans{300}
\loigiai{
\textbf{Đáp số:} $300$\\
$15 \times 20 = 300\text{ (cm}^2\text{)}$.
}
\end{ex}

\begin{ex} %[5H5H28]
Cắt tấm bìa hình vuông cạnh $20\text{ cm}$ thành $4$ hình tam giác bằng nhau. Tính diện tích mỗi tam giác theo xăng-ti-mét vuông.
\shortans{100}
\loigiai{
\textbf{Đáp số:} $100$\\
$20 \times 20 : 4 = 100\text{ (cm}^2\text{)}$.
}
\end{ex}

\begin{ex} %[5H5V28]
Cắt $4$ góc hình vuông cạnh $5\text{ cm}$ từ tấm bìa hình vuông cạnh $30\text{ cm}$. Tính diện tích phần bìa còn lại theo xăng-ti-mét vuông.
\shortans{800}
\loigiai{
\textbf{Đáp số:} $800$\\
$30 \times 30 - 4 \times 5 \times 5 = 900 - 100 = 800\text{ (cm}^2\text{)}$.
}
\end{ex}

\begin{ex} %[5H5V28]
Tính diện tích phần nằm giữa hình vuông cạnh $10\text{ cm}$ và hình tròn có đường kính $10\text{ cm}$ nằm bên trong hình vuông theo xăng-ti-mét vuông. (Nhập kết quả dưới dạng số thập phân).
\shortans{21,5}
\loigiai{
\textbf{Đáp số:} $21{,}5$\\
$10 \times 10 - 3{,}14 \times 5 \times 5 = 100 - 78{,}5 = 21{,}5\text{ (cm}^2\text{)}$.
}
\end{ex}

\begin{ex} %[5H5V28]
Một cánh diều hình thoi có hai đường chéo dài $60\text{ cm}$ và $80\text{ cm}$. Tính diện tích cánh diều theo xăng-ti-mét vuông.
\shortans{2400}
\loigiai{
\textbf{Đáp số:} $2400$\\
$\dfrac{60 \times 80}{2} = 2\,400\text{ (cm}^2\text{)}$.
}
\end{ex}
